\section{ドップラー効果(2)・屈折の法則}

\subsection{斜めドップラー効果}
ここまでは波源と観測者が，その2点を結ぶ直線方向のみに運動する場合を考えたが，それ以外の方向に波源・観測者が運動する場合は，波源と観測者の速度のうち，波の伝播方向の速度成分のみがドップラー効果に寄与する．

\begin{wrapfigure}[8]{r}[5mm]{90mm}
  \centering
  \vspace{-5mm}
  \includegraphics[keepaspectratio,width=90mm]
  {text_p113_fig1.png}
\end{wrapfigure}
例えば，右図のように波源と観測者が運動する場合は，波の伝播方向の速度成分である$v\cos \theta$[m/s]および$u\cos \varphi$[m/s]がドップラー効果に有効な速度成分であるから，
\[f'=\dfrac{c+u\cos\varphi}{c-v\cos \theta} f \hspace{2mm}[{\rm Hz}]\]
となる振動数の音が観測される．

\subsection{衝撃波}
今まで考えてきたドップラー効果はいずれも波源の速さが波の速さを超えない場合であった．
もし波源の速さが波の速さを超えるようになると，波源の前方へ進む波の波長はドップラー効果の式から計算すると負になってしまう．
すなわち，{\bf 波源の速さが波の速さを超えると波源の前方ではドップラー効果は起こらない．}

\begin{wrapfigure}[18]{r}[5mm]{65mm}
  \centering
  \vspace{-5mm}
  \includegraphics[keepaspectratio,width=65mm]
  {text_p115_fig1.png}
\end{wrapfigure}
音源Sの速さが$v$[m/s]，音速が$c$[m/s]とし，$v>c$の条件が満たされている場合について考えると，図$\fbox{1}$のP$_1$を通過したときに出た音波はP$_1$を中心として同心円上に広がる．
やがてP$_2$に音源は達するが，このときのP$_1$を通過したときに出された音波の波面は図$\fbox{2}$のようになっている．
P$_2$を通過するときに出す音波は通過したときからP$_2$を中心として同心円上に広がっていく．
これを繰り返していくと，各P$_1$，P$_2$，P$_3$などから出た音波は図のような円錐形の波面を作ることになる．
波面上には各時刻に出された球面波が重なり合うため，非常に大きな振幅を持つ波面となる．
このような波を{\bf 衝撃波}という．

P$_1$を通過してから時間$t$[s]だけ経過したときを考えると，衝撃波の半頂角を$\theta$として，図より
\[\sin \theta =\dfrac{c\cdot t}{v\cdot t}=\dfrac{c}{v}\]
という関係が満たされる．
衝撃波においては，波源の進行方向の波源より前方には波は伝わっていかず，波源の後方へ進む波については通常のドップラー効果が観測される．

\vspace{5mm}
 \begin{matome}{衝撃波}
 \hspace{3mm} 波源の速さが波の速さを超えると，波源の進行方向の波源より前方には波は伝わっていかない．\\
 \hspace{3mm} 波源から出る波の波面は波源を頂点とする円錐形をなし，その半頂角$\theta$は$\sin \theta=\dfrac{c}{v}$を満たす．
  \end{matome}
 
\subsection{光速度}
光は真空中において{\bf 高速度}\mbox{\boldmath $c=2.99792458\times 10^8 \fallingdotseq 3.00 \times 10^8$} {\bf [m/s]の速さをもつ横波}であることが知られている．
そのため，波動としての特徴である回折や干渉を起こし，また横波であるために後に述べる偏光といった現象を起こす．
真空中の{\bf 光の速さは光の波長，振動数によらず一定}であることが知られている．

 \begin{matome}{光速度}
 \hspace{3mm} $c\fallingdotseq 3.00\times 10^8$[m/s]\hspace{5mm} 波長・振動数によらない
 \end{matome}

\subsection{光の屈折}
ある媒質を伝わる光波が別の媒質に入射するとき，光波はその一部が境界面で反射し，残りは屈折してもう一方の媒質に入る．
このとき，境界面に対して垂直に立てた法線に対する角度をそれぞれ{\bf 入射角，屈折角}という．

\begin{wrapfigure}[8]{r}[5mm]{60mm}
  \centering
  \vspace{-5mm}
  \includegraphics[keepaspectratio,width=60mm]
  {text_p116_fig1.png}
\end{wrapfigure}
入射角を$i$，屈折角を$r$とし，媒質I，媒質II中の光の速さ，波長をそれぞれ$c_1,c_2$[m/s]，$\lambda_1,\lambda_2$[m]とすると，
\[\mbox{\boldmath $\dfrac{\sin i}{\sin r} =\dfrac{c_1}{c_2}= \dfrac{\lambda_1}{\lambda_2} =n_{12} $}\hspace{5mm} \cdots \cdots （☆）\]
が成立する．
この比の値は入射角$i$に依存しない．
$n_{12}$を{\bf 媒質Iに対する媒質IIの屈折率（相対屈折率）}という．
また，{\bf 屈折の前後で振動数は変化しない．}
この式は光波に限らず，波動全般について成立する．

特に光波の場合は，真空に対する媒質の屈折率をその媒質の{\bf 絶対屈折率}という．
媒質I，媒質IIの絶対屈折率を$n_1,n_2$，真空中の光の速さを$c$[m/s]とすると，屈折率の定義式より
\[ n_1=\dfrac{c}{c_1},\hspace{5mm}n_2=\dfrac{c}{c_2}\]
が成立する．
$n_{12},n_1,n_2$の間には$n_{12}=\dfrac{n_2}{n_1}$が成立する．
絶対屈折率$n$は{\bf 真空の場合は1，それ以外は1より大きくなる}と考えてよい．

なお，$n_1,n_2$を用いると，☆式は次の形に書き直すことができる．
\[\mbox{\boldmath $ n_1\sin i=n_2\sin r ,\hspace{5mm} n_1c_1=n_2c_2 ,\hspace{5mm} n_1\lambda_1=n_2\lambda_2$}\]

 \begin{matome}{屈折の法則}
 \hspace{3mm} $n_1\sin i=n_2\sin r ,\hspace{5mm} n_1c_1=n_2c_2 ,\hspace{5mm} n_1\lambda_1=n_2\lambda_2$\\
 屈折の前後で波の振動数は不変．
 \end{matome}

屈折率の大きなものから小さなものへ光が進む場合は入射角$i$よりも屈折角$r$の方が大きくなるため，入射角がある値$\theta_0$を越えると屈折光は存在できなくなり，すべての光が反射される．
これを全反射といい，角度$\theta_0$を{\bf 臨界角}という．

\begin{wrapfigure}[6]{r}[5mm]{60mm}
  \centering
  \vspace{-5mm}
  \includegraphics[keepaspectratio,width=60mm]
  {text_p117_fig1.png}
\end{wrapfigure}
簡単のため，右図のように屈折率$n$の媒質中の点光源から出た光が真空中へと出ていく場合について考える．
屈折の法則から$\sin r=n\sin i$となるが，\mbox{\boldmath $\sin r=n\sin i >1$}{\bf となると屈折光は存在できない．}
すなわち，
\[\mbox{\boldmath $ \sin i>\dfrac{1}{n}$}\]
なる入射角$i$の光は全反射することになる．
よって，臨界角を$\theta_0$とすると，
\[\sin \theta_0 =\dfrac{1}{n}\]
が成立する．
{\bf 臨界角を求めるには，屈折の法則において屈折角を直角と置けばよい．}

\subsection{ホイヘンスの原理と屈折の法則}
波が平面上もしくは空間中を伝わっていくとき，ある瞬間に波の同位相の点を連ねていくと，一つの曲線もしくは曲面ができる．
この面を波面といい，この波面が平面である波を{\bf 平面波}，球面になる波を{\bf 球面波}という．
波面は常に波の進行方向に対して垂直になっている．

\begin{wrapfigure}[13]{r}[5mm]{100mm}
  \centering
  \vspace{-5mm}
  \includegraphics[keepaspectratio,width=100mm]
  {text_p118_fig1.png}
\end{wrapfigure}
波面が次の瞬間にどのように広がっていくかは明らかなことも多いが，一般的には次に述べる{\bf ホイヘンスの原理}に従うことが経験的に知られている．
波が届いた媒質の点は波と同じ振動数で振動を始める．
そのため，これらの点は新たな波源とみなすことができる．
ある波面が次の瞬間にどのような波面になるかは，波面上の各点が新たな波源となって出す球面波（これを{\bf 素元波}という）がどのような曲面を形成するかを考えればよい．
多数の素元波のすべてに接する曲面（包絡面）のうち，波の進行方向側にできるものが次の瞬間の波面となる．
これを{\bf ホイヘンスの原理}といい，これを用いると波の反射，屈折，回折現象を説明することができる．

 \begin{matome}{ホイヘンスの原理}
 \hspace{3mm} 波動がある瞬間に作る波面の各点から出る無数の素元波の包絡面が次の瞬間の波面となる．
  \end{matome}
 
 \begin{wrapfigure}[7]{r}[5mm]{60mm}
  \centering
  \vspace{-5mm}
  \includegraphics[keepaspectratio,width=60mm]
  {text_p118_fig2.png}
\end{wrapfigure}
 波が障害物に当たったとき，幾何学的には影となる部分にも実際には若干の波が回り込むことが知られており，これを{\bf 波の回折}といい，回り込んだ波を{\bf 回折波}という．
 
 回折が顕著に起こるか否かは障害物やすき間の大きさと波の波長の大小関係によって決まる．
 すなわち，{\bf すき間などの大きさに比べて十分に大きい波長を持つ波は回折が顕著であり，影の部分に回り込む程度が大きい．}
 逆に，すき間などの大きさに比べて小さい波長を持つ波は回折しにくいことが知られており，これはホイヘンスの原理によって説明できることが知られている．
 
  \begin{figure}[h]
\begin{center}
\includegraphics[width=110mm]{text_p118_fig3.png}
\end{center}
\end{figure}
  \vspace{40mm}
  \begin{matome}{回折現象}
 \hspace{3mm} 波の波長が障害物やすき間の大きさに比べて同程度より大きくなると，波の回折が顕著となり，波は幾何学的に影となる部分によく回り込むようになる．
  \end{matome}
  
  
\begin{wrapfigure}[10]{r}[5mm]{90mm}
  \centering
  \vspace{-5mm}
  \includegraphics[keepaspectratio,width=90mm]
  {text_p119_fig1.png}
\end{wrapfigure}
以上で述べたホイヘンスの原理を用いると，屈折の法則の公式を以下のようにして導くことができる．
図のように，媒質Iから媒質IIへ向けて位相の揃った平面波が入射角$i$で入射し，屈折角$r$で屈折していく場合を考えよう．
波面は波の進行方向に対して垂直であるから，図のA,A$^\prime$は同一波面上に，また屈折した後のB,B$^\prime$も同一波面上にある．
すなわち，A,A$^\prime$は同位相であり，B,B$^\prime$も同位相であるから，(a)のルートを通る波がAからBへ進む間に，(b)のルートを通る波はA$^\prime$からB$^\prime$へ進んでいるはずである．

ここで，各媒質中の波の速さおよび波長をそれぞれ$c_1,c_2$[m/s]，$\lambda_1,\lambda_2$[m]とする．
AからBへ波が進むのにかかる時間（=A$^\prime$からB$^\prime$へ波が進むのにかかる時間）を$T$[s]とすると，波の進行距離および幾何学的関係より，
\[\overline{{\rm A}^\prime {\rm B}^\prime} =c_1T,\hspace{4mm} \overline{{\rm AB} } =c_2T,\hspace{4mm} \overline{{\rm A}^\prime {\rm B}^\prime} =\overline{{\rm AB}^\prime}\sin i ,\hspace{4mm} \overline{{\rm AB}} = \overline{ {\rm AB}^\prime}\sin r\]
である．
これらの式より，
\[\dfrac{\overline{{\rm A}^\prime {\rm B}^\prime}}{\overline{{\rm AB} }} =\dfrac{\overline{{\rm AB}^\prime}\sin i }{ \overline{ {\rm AB}^\prime}\sin r} =\dfrac{c_1T}{c_2T}\]
となるので，
\[\dfrac{\sin i}{\sin r}=\dfrac{c_1}{c_2}\]
なる屈折の法則が成立する．
また，屈折の前後で波の振動数$f$は不変であるから，
\[ \dfrac{\sin i}{\sin r}=\dfrac{c_1}{c_2}=\dfrac{c_1/f}{c_2/f}=\dfrac{\lambda_1}{\lambda_2}\]
となる．

\newpage

\begin{figure}[h]
\begin{center}
\includegraphics[width=170mm]{prob_p61.png}
\vspace{150mm}
\end{center}
\end{figure}

\newpage

\begin{figure}[h]
\begin{center}
\includegraphics[width=170mm]{prob_p62.png}
\vspace{150mm}
\end{center}
\end{figure}

\newpage

\begin{figure}[h]
\begin{center}
\includegraphics[width=170mm]{prob_p63.png}
\vspace{150mm}
\end{center}
\end{figure}

\newpage

\begin{figure}[h]
\begin{center}
\includegraphics[width=170mm]{prob_p171.png}
\vspace{150mm}
\end{center}
\end{figure}

\newpage

\begin{figure}[h]
\begin{center}
\includegraphics[width=170mm]{prob_p172.png}
\vspace{150mm}
\end{center}
\end{figure}

\newpage

\begin{figure}[h]
\begin{center}
\includegraphics[width=170mm]{prob_p173.png}
\vspace{150mm}
\end{center}
\end{figure}

\newpage

\begin{figure}[h]
\begin{center}
\includegraphics[width=170mm]{prob_p174.png}
\vspace{150mm}
\end{center}
\end{figure}

\newpage

\begin{figure}[h]
\begin{center}
\includegraphics[width=170mm]{prob_p175.png}
\vspace{150mm}
\end{center}
\end{figure}
